% !TEX root = ../main.tex
\section*{2. Related Work}
Language models have been studied in social dilemmas, always as equals. They play repeated two-player games competently but unforgivingly \citep{akata2025repeated}, and most groups sharing a resource collapse unless agents reason about the effect of their actions on everyone \citep{piatti2024govsim}. In these worlds agents have ranks, and each one's survival depends on those above and below it.

Simulation frameworks come closer in design. Concordia's Game Master, itself a language model, turns actions stated in words into their effects \citep{vezhnevets2023concordia}. My question required the reverse: identical rules in every world, so that a difference between two worlds can be traced to what I changed. In place of the Game Master I use an engine that applies the rules, records every prompt and reply, and can check a world by rerunning it from those records.

Groups of agents fail through miscoordination, conflict and collusion \citep{hammond2025risks}, malicious instructions can pass from agent to agent \citep{lee2024infection}, and agents can collude in secret \citep{motwani2024collusion}. \citet{dafoe2020open} pose the open problems of cooperative AI without running experiments, and name four capabilities it requires: understanding, communication, commitment and institutions. Each appears in these worlds as something I can change: what agents may see, the messages they may send, the promises that bind them, and the hierarchy and its punishments.
